
\documentclass{article}
\usepackage{amsmath}
\usepackage{enumerate}
\begin{document}

\documentclass{article}
\usepackage{amsmath}
\usepackage{enumerate}
\begin{document}

\documentclass{article}
\usepackage{amsmath}
\usepackage{enumerate}
\begin{document}
### Chapter 2: Thermodynamic Properties and Molecular Dynamics

#### 2.1 Pressure from Molecular Collisions

**Molecular Collision Dynamics**

- **D/ = ` ∙∆_ ∙**: This equation represents the relationship between molecular collisions and the dynamic process of gas molecules interacting with a surface or wall.

- **3. Momentum Transferred to the Wall**
   - **h< = 2i ∙E<**: The momentum transferred per collision is given by this formula, where `2i` refers to mass and `E<` represents velocity in a specific direction.
   
- **4. Total Transferred Momentum**
   - **j = i ∙`50 = o50** 
     - j ∙` = 5 ∙0**: This equation describes the total pressure generated by molecular collisions, where `i`, `j`, and `o` are mass-related constants and variables.

#### 2.2 The Ideal Gas Assumption

- **Ideal gas**: An ideal gas is a theoretical model of gases in which molecules have no intermolecular interactions or attraction/repulsion forces.
   
- **Elastic Collisions**:
   - Elastic collisions between molecules and the wall are not necessary if the system is at thermal equilibrium, meaning `8 666668`.

#### 2.3 Equipartition of Energy

The energy distribution in an ideal gas can be analyzed using statistical mechanics:

1. **Translational Kinetic Energy**:
   - Each degree of freedom contributes ∙4 ∙0 to the average translational kinetic energy.
   
2. **Rotational Degrees of Freedom**:
   - Rotational degrees are excited if atoms' centers of gravity are not aligned with the rotation axis.

3. **Vibrational Degrees of Freedom**:
   - Vibrational energy per degree of freedom is 4 ∙0, where quanta of energy `ℎ∙` must be considered for quantum effects at low temperatures.

#### Summary Table: Degrees of Freedom

|**Number of Atoms**|yK|yL|yM|**Example**|
|---|---|---|---|---|
|1|3|0|0|Ar|
|2|3|2|1|N2|
|3,linear|3|2|4|O=C=O|
|3, non-linear|3|3|3|H2O|
|_N,_linear|3|2|3_N_-5|H-CºC-H|
|_N,_non-linear|3|3|3_N_-6|C2H5|

#### 2.4 Thermodynamic Properties and Degrees of Freedom

- **Degrees of freedom (s%)**: The sum includes translational, rotational, and vibrational contributions.
  
- **Energy per molecule**:
   - Translation: ∙4 ∙0 
   - Rotation: ∙4 ∙0
   - Vibration: 2 ∙4 ∙0

#### Remarks 

1. **Boltzmann's Constant**: Understanding the relationship between Boltzmann’s constant and the universal gas constant.
   
2. **Phase Transition Dynamics**:
   - The behavior of isentropic expansion in the two-phase area liquid-vapor can have a positive inclination if conditions meet certain criteria.

#### Key Equations

- **Momentum Transfer**: h< = 2i ∙E<
  
- **Total Momentum and Pressure**: 
   - j = i ∙`50
   - j ∙` = 5 ∙0 

- **Energy Contribution from Degrees of Freedom**:
   - r = (s% + s- + 2 ∙s.)<sup>!</sup> ∙4 ∙0 8

#### Examples for Calculation at 300K
- Specific examples and calculations would be included here, considering the energy contributions based on degrees of freedom.

**Notes:**
- Understanding the interplay between molecular dynamics (collisions, rotation, vibration) and thermodynamic properties (pressure, temperature) is crucial.
- The role of quantum mechanics in determining vibrational excitation at low temperatures needs careful consideration.\end{document}